\begin{examplebox}
	\textbf{\textcolor{CExample}{例 1（基础）}}

	\textbf{\textcolor{CExample}{题目：}}
	求函数 $y=3\sin(2x+\dfrac{\pi}{3})$ 在 $[0,\dfrac{\pi}{2}]$ 上的最大值和最小值。

	\textbf{\textcolor{CExample}{解题步骤：}}
	\begin{description}[style=nextline, leftmargin=2.8em, labelsep=0.5em, itemsep=0.45em]
		\item[\textbf{第一步（换元定范围）：}] 令 $t=2x+\dfrac{\pi}{3}$，由 $x\in[0,\dfrac{\pi}{2}]$ 得 $t\in[\dfrac{\pi}{3},\dfrac{4\pi}{3}]$。
			\par\textbf{理由：} 线性映射，端点对应。

		\item[\textbf{第二步（正弦最值）：}] 在 $[\dfrac{\pi}{3},\dfrac{4\pi}{3}]$ 上，$\sin t$ 在 $t=\dfrac{\pi}{2}$ 处取得最大值 $1$，在 $t=\dfrac{4\pi}{3}$ 处取得最小值 $-\dfrac{\sqrt{3}}{2}$（因为 $\dfrac{3\pi}{2}$ 不在区间内，且 $\sin\dfrac{4\pi}{3}<\sin\dfrac{\pi}{3}$）。
			\par\textbf{理由：} 正弦函数图象或单位圆。

		\item[\textbf{第三步（乘系数得结果）：}] $y_{\max}=3\times1=3$，$y_{\min}=3\times(-\dfrac{\sqrt{3}}{2})=-\dfrac{3\sqrt{3}}{2}$。
			\par\textbf{理由：} 由结论 $y_{\max}=A\cdot M$，$y_{\min}=A\cdot m$。
	\end{description}

	\textbf{\textcolor{CExample}{关键结论：}}
	\textbf{最大值为 $3$，最小值为 $-\dfrac{3\sqrt{3}}{2}$。}

	\medskip
	\textbf{\textcolor{CExample}{例 2（稍变形）}}

	\textbf{\textcolor{CExample}{题目：}}
	求函数 $y=2\sin(\dfrac{x}{2}-\dfrac{\pi}{4})+1$ 在 $[0,4\pi]$ 上的值域。

	\textbf{\textcolor{CExample}{解题步骤：}}
	\begin{description}[style=nextline, leftmargin=2.8em, labelsep=0.5em, itemsep=0.45em]
		\item[\textbf{第一步（换元定范围）：}] 令 $t=\dfrac{x}{2}-\dfrac{\pi}{4}$，由 $x\in[0,4\pi]$ 得 $t\in[-\dfrac{\pi}{4},\dfrac{7\pi}{4}]$。
			\par\textbf{理由：} 线性映射。

		\item[\textbf{第二步（正弦最值）：}] 在 $t\in[-\dfrac{\pi}{4},\dfrac{7\pi}{4}]$ 上，$\sin t$ 的最大值为 $1$（在 $t=\dfrac{\pi}{2}$ 取得），最小值为 $-1$（在 $t=\dfrac{3\pi}{2}$ 取得）。
			\par\textbf{理由：} 利用正弦图象，区间包含最大值点和最小值点。

		\item[\textbf{第三步（加常数得结果）：}] $y_{\max}=2\times1+1=3$，$y_{\min}=2\times(-1)+1=-1$。
			\par\textbf{理由：} 先由结论得正弦部分最值，再加常数。
	\end{description}

	\textbf{\textcolor{CExample}{关键结论：}}
	\textbf{值域为 $[-1,3]$。}

	\medskip
	\textbf{\textcolor{CExample}{例 3（综合应用）}}

	\textbf{\textcolor{CExample}{题目：}}
	求函数 $y=4\sin(3x+\dfrac{\pi}{6})$ 在 $[-\dfrac{\pi}{18},\dfrac{\pi}{9}]$ 上的最大值和最小值。

	\textbf{\textcolor{CExample}{解题步骤：}}
	\begin{description}[style=nextline, leftmargin=2.8em, labelsep=0.5em, itemsep=0.45em]
		\item[\textbf{第一步（换元映射）：}] 令 $t=3x+\dfrac{\pi}{6}$，由 $x\in[-\dfrac{\pi}{18},\dfrac{\pi}{9}]$ 得 $t\in[0,\dfrac{\pi}{2}]$。
			\par\textbf{理由：} 计算端点：$t_1=3\times(-\dfrac{\pi}{18})+\dfrac{\pi}{6}=0$，$t_2=3\times\dfrac{\pi}{9}+\dfrac{\pi}{6}=\dfrac{\pi}{2}$。

		\item[\textbf{第二步（正弦最值）：}] 在 $[0,\dfrac{\pi}{2}]$ 上，$\sin t$ 单调递增，最大值为 $\sin\dfrac{\pi}{2}=1$，最小值为 $\sin0=0$。
			\par\textbf{理由：} 正弦函数在 $[0,\dfrac{\pi}{2}]$ 上单调递增。

		\item[\textbf{第三步（乘系数得结果）：}] $y_{\max}=4\times1=4$，$y_{\min}=4\times0=0$。
			\par\textbf{理由：} 由结论 $y_{\max}=A\cdot M$，$y_{\min}=A\cdot m$。
	\end{description}

	\textbf{\textcolor{CExample}{关键结论：}}
	\textbf{最大值为 $4$，最小值为 $0$。}
\end{examplebox}
